% Generated by tools/edn_latex.py from reports/edn.md.
% Do not edit: change reports/edn.md and run `make edn`.
\ednentry{
  id = {EDN-22},
  title = {Drop listings registered after the age reference date},
  note = {\textbf{Evidence corrected on 2026-09-29.} The decision said preprocessing drops ``the 164 listings''. 164 is the raw-file count;
preprocessing runs after scoping, so it removes 27 of them. The alternatives below already said 26, measured after deduplication.
The decision itself is unchanged.
\par
\textbf{Description corrected on 2026-09-30.} The decision described the 27 as what survives ``the EDN-04 scope and the training price range''.
The rule order implemented in issue \#34 puts this rule above the price range, so the 27 survive the scope alone.
The number is the same either way, because all 27 are inside the range; only the description of how it was obtained was wrong.
The decision itself is unchanged.},
  date = {2026-09-29},
  milestone = {M3: Quality Assurance},
  topic = {Data Validation and Preprocessing},
  participants = {@lukas2510},
  decision = {Preprocessing drops every listing whose \texttt{registration\_\allowbreak{}date} is after the age reference date (the 2025-11-08 snapshot date in training, the request date in serving). There are 164 such listings in the raw file, of which \textbf{27} survive the EDN-04 scope, the one rule the implemented order puts above this one, and so are the ones preprocessing actually removes; all 27 are inside the training price range, which runs after this rule, so the count is the same whichever of the two goes first; 26 remain if deduplication runs first, the difference being one duplicate row. The Great Expectations suite asserts \texttt{registration\_\allowbreak{}date <=\allowbreak{} reference date} on the processed data as a hard expectation, and the same rule on the raw data with \texttt{mostly=\allowbreak{}0.\allowbreak{}99}, so a future scrape that suddenly carries many such rows is flagged instead of silently cleaned away.},
  rationale = {\emph{Alternatives considered:}
\begin{itemize}[leftmargin=*, topsep=2pt]
\item \textbf{Option A (chosen): drop the rows.}
\newline Pros: 164 of 118,382 raw rows is 0.14\,\%, and only 26 of them survive the EDN-04 scope, so nothing measurable is lost; makes \texttt{age >=\allowbreak{} 0} a genuine invariant that the pipeline, the tests and the API contract can all rely on; no special case anywhere in the feature code.
\newline Cons: throws away the other, possibly correct attributes of those listings; the rule has to live in preprocessing, so the raw expectation has to tolerate what preprocessing removes.
\item \textbf{Option B: clamp the age at zero.}
\newline Pros: keeps the rows and their remaining attributes; a single \texttt{max(0, .\allowbreak{}.\allowbreak{}.\allowbreak{})} in the feature code.
\newline Cons: invents a registration date that the source does not support, and does so for rows whose date is demonstrably wrong (the latest is 2026-11-01, almost a year after the snapshot); a clamped zero is indistinguishable from a genuinely new car, which is exactly the group EDN-04 scopes out.
\item \textbf{Option C: keep the rows with a negative age.}
\newline Pros: no data loss, no cleaning rule; gradient boosting tolerates the values.
\newline Cons: the API would have to accept a negative age too, or training and serving would disagree; removes the cheapest available sanity check on the most important feature.
\item \textbf{Option D: repair the date from \texttt{production\_\allowbreak{}year}.}
\newline Pros: would keep the rows with a defensible value.
\newline Cons: \texttt{production\_\allowbreak{}year} is only 19\,\% filled, so it repairs almost none of them; a repair rule needs its own validation for a group this small.
\end{itemize}
\par
\emph{Why:} For the rows this decision actually governs, a registration date after the snapshot is a data-entry error rather than a rare but real phenomenon: 137 of the 164 sit in a single month, January 2026, which looks like a year typo. The reading is weaker for the 114 new cars among them, where a dealer entering a planned first-registration date is plausible, but those fall outside the used-car scope before preprocessing runs; only 26 of the 164 reach it. At that share the cost of dropping them is nil, while every alternative either fabricates a value or gives up the \texttt{age >=\allowbreak{} 0} invariant. Bounding the raw data with \texttt{mostly=\allowbreak{}0.\allowbreak{}99} rather than a hard rule keeps the raw suite honest about what the source actually contains while still catching a systematic worsening, which matters because the \texttt{ES} holdout and the M6 drift scenario re-run these checks on data we have not profiled.},
  ai = {Information seeking, Alternative generation, Alternative assessment, Recommendation},
  response = {Accepted},
  assessment = {AI profiled the 164 rows, laid out the four options above and recommended dropping them, together with the split between a hard expectation on the processed data and a \texttt{mostly} bound on the raw data. Lukas accepted the recommendation because the dropped share is negligible and the invariant is worth more than the rows.},
  aievidence = {Claude Code session on 2026-09-29 while analysing \href{https://github.com/mlops-2627q1-mds-upc/MLOPS_Recommenditos/issues/25}{issue \#25}: AI presented the options as a decision question, Lukas chose ``drop them''.},
  evidence = {\href{https://github.com/mlops-2627q1-mds-upc/MLOPS_Recommenditos/issues/25}{Issue \#25}, \href{https://github.com/mlops-2627q1-mds-upc/MLOPS_Recommenditos/pull/24}{PR \#24}, \href{https://github.com/mlops-2627q1-mds-upc/MLOPS_Recommenditos/blob/main/docs/docs/problem-spec.md}{problem specification} section 2.},
}
